\section{Introduction}

Public verification has a sharp privacy edge. A voter should be able to ask
whether a counted or accepted record is present without gaining an artifact that
can prove a candidate choice to a buyer, employer, family member, campaign, or
state actor. The hard rule for RCOUNT is therefore simple: prove inclusion,
never prove choices.

V.3 described how source and package hashes make public evidence tamper-evident.
V.4 narrows the problem to voter-facing proof material. The baseline proof is
not a complete cryptographic voting system. It is a privacy gate over public
proof records: if the record carries candidate selections or linkable voter
facts, it fails.

The voter-facing interpretation is deliberately narrow. A passing proof can
support the statement that an anonymized accepted-token or counted-token record
appears in the public package. It does not support the stronger statement that
``my vote for candidate X was counted,'' because that stronger statement is
exactly the receipt a coercer would want. Refusing to publish that evidence is a
voter protection, not a missing feature.
